The fourth-order Runge-Kutta method is a widely used numerical technique that provides excellent accuracy while maintaining reasonable computational complexity.
Again, considering the same initial value problem as used in ~\ref{subsubsec:euler-method}, the method requires four slope evaluations per step:

\[
\begin{array}{ll}
k_1 = f(t_n, y_n) &
k_3 = f\left(t_n + \frac{\Delta t}{2}, y_n + \frac{\Delta t}{2}k_2\right) \\[4pt]
k_2 = f\left(t_n + \frac{\Delta t}{2}, y_n + \frac{\Delta t}{2}k_1\right) &
k_4 = f\left(t_n + \Delta t, y_n + \Delta t \cdot k_3\right)
\end{array}
\]

These four slopes are then combined in a weighted average to compute the next value:

\[
y_{n+1} = y_n + \frac{\Delta t}{6}(k_1 + 2k_2 + 2k_3 + k_4).
\]

Despite requiring more function evaluations per step, its higher order of accuracy often allows for larger step sizes while maintaining precision.
The fourth-order Runge-Kutta method strikes an excellent balance between computational cost and accuracy, making it one of the most popular choices for solving ordinary differential equations.